
\documentclass[10pt,
% a4paper,
%twocolumn,
fleqn,
%landscape, 
% papersize,
dvipdfmx,
uplatex
]{jsarticle}



\def\maru#1{\textcircled{\scriptsize#1}}%丸囲み番号

% \RequirePackage[2020/09/30]{platexrelease}

%太字設定
\usepackage[deluxe]{otf}

\usepackage{emathEy}

\usepackage[g]{esvect}

%大きな文字
\usepackage{fix-cm}

%定理環境
\usepackage{emathThm}
%\theoremstyle{boxed}
\theorembodyfont{\normalfont}
\newtheorem{Question}{問題}[subsection]
\newtheorem{Q}{}[subsection]
\newtheorem{question}[Question]{}
\newtheorem{quuestion}{}[subsection]

%セクション，大問番号のデザイン
\renewcommand{\labelenumi}{(\arabic{enumi})}
\renewcommand{\theenumii}{\alph{enumii})}
\renewcommand{\thesection}{第\arabic{section}章}

%用紙サイズの詳細設定
\usepackage{bxpapersize}
\papersizesetup{size={80mm,45mm}}
\usepackage[top=0.7zw,bottom=0truemm,left=3truemm,right=133truemm]{geometry}
\usepackage[dvipdfmx]{graphicx}

%余白など
\usepackage{setspace} % 行間
\setlength{\mathindent}{1zw}
\setlength\parindent{0pt}


%色カラーに関する設定
\usepackage{color}
\definecolor{shiro}{rgb}{0.95703125,0.87109375,0.7421875}
\definecolor{kin}{rgb}{0.95703125,0.87109375,0.7421875}
\definecolor{orange}{rgb}{1,0.7,0.2}
\definecolor{bradorange}{rgb}{1,0.5,0}
\definecolor{pink}{rgb}{0.9176,0.5686,0.5960}
\definecolor{mizu}{rgb}{0.6156,0.8,0.9955}
\color{kin}
% \pagecolor{hukamido}

\usepackage{at}%図の配置
% \usepackage{wallpaper}

\begin{document}

\bf\boldmath

\at(0cm,0cm){\includegraphics[width=8cm,bb=0 0 1920 1080]{./youtube/thumbnails/templates/smart_background/図形と方程式.jpeg}}
\at(6.6cm,0.2cm){\small\color{bradorange}$\overset{\text{図形と方程式}}{\text{応用}}$}
{\color{orange}\Large\underline{片側が固定された線分の中点の軌跡}}\vspace{0.3zw}

\Large 
問．円$x^2+y^2=1$上の動点$\text{P}$と，\\
点$\text{A}\left(3,\;4\right)$を結ぶ

\huge
\hspace{0.5zw}線分APの中点$\text{M}$の\vspace{-0.2zw}\\
\hfill 軌跡を求めよ．

\newpage



\at(0cm,0cm){\includegraphics[width=8cm,bb=0 0 1920 1080]{./youtube/thumbnails/templates/smart_background/図形と方程式.jpeg}}
{\color{orange}\bf\boldmath\huge\underline{線形計画法}}\vspace{0.3zw}

\normalsize
\bf\boldmath 問．

\Large
\vspace{-1zw}
\hspace{1zw}
$\left\{\begin{array}{l}3x+y\geqq 6\;\\x+3y\geqq 6\;\\x+y\leqq 6\end{array}\right.$のとき，

\Large
\vspace{0.5zw}
\hspace{0.5zw}
$x+2y$の最大値と最小値\large を求めよ．

\at(6.6cm,0.2cm){\small\color{bradorange}$\overset{\text{図形と方程式}}{\text{典型}}$}

\newpage

\at(0cm,0cm){\includegraphics[width=8cm,bb=0 0 1920 1080]{./youtube/thumbnails/templates/smart_background/図形と方程式.jpeg}}
{\color{orange}\bf\boldmath\huge\underline{領域の図示 全盛り}}\vspace{0.3zw}

\small
\bf\boldmath 問．次の不等式が表す領域を図示せよ．

\small
(1)  $y>x^2+1$ \hspace{2.1zw} (2)  $3x-2y-2\geqq 0$\vspace{0.3zw}\\
(3)  $x\leqq 2$  \hspace{5zw} (4)  $\left(x+2\right)^2+y^2<1$\vspace{0.3zw}\\
(5)  $x^2+y^2-6x-2y+1\geqq 0$\vspace{0.3zw}\\
(6)  $\left\{\begin{array}{l}\hspace{-0.5zw}x^2+y^2<{25}\\\hspace{-0.5zw}y<3x-5\end{array}\right.$
\hspace{-0.4zw}(7)  $\left(x-y\right)\left(x+y-2\right)<0$\vspace{0.3zw}\\

\at(6.6cm,0.2cm){\small\color{bradorange}$\overset{\text{図形と方程式}}{\text{典型}}$}

\newpage

\at(0cm,0cm){\includegraphics[width=8cm,bb=0 0 1920 1080]{./youtube/thumbnails/templates/smart_background/図形と方程式.jpeg}}
{\color{orange}\bf\boldmath\LARGE\underline{三角形の重心の軌跡}}\vspace{0.3zw}

\normalsize
\bf\boldmath 問．$2$点$\text{O}\left(0,\;0\right),\;\text{A}\left(1,\;0\right)$と\\
\hfill 円$x^2+y^2=9$上を動く点$\text{Q}$を頂点とする

\huge
\hspace{0.0zw}三角形$\text{OAQ}$の重心$\text{P}$\\
\hfill の軌跡を求めよ.
\at(6.6cm,0.2cm){\small\color{bradorange}$\overset{\text{図形と方程式}}{\text{典型}}$}

\newpage

\at(0cm,0cm){\includegraphics[width=8cm,bb=0 0 1920 1080]{./youtube/thumbnails/templates/smart_background/図形と方程式.jpeg}}
{\color{orange}\bf\boldmath\huge\underline{軌跡 Lv.1}}\vspace{0.5zw}

\large 
\bf\boldmath 問．$2$点$\text{A}\left(-3,\;0\right),\;\text{B}\left(2,\;0\right)$からの

\huge
\vspace{0.2zw}
\hspace{0.0zw}
距離の比が$3:2$である\\
\hfill 点$\text{P}$の軌跡を求めよ．

\at(6.6cm,0.2cm){\small\color{bradorange}$\overset{\text{図形と方程式}}{\text{典型}}$}

\newpage

\at(0cm,0cm){\includegraphics[width=8cm,bb=0 0 1920 1080]{./youtube/thumbnails/templates/smart_background/図形と方程式.jpeg}}
{\color{orange}\bf\boldmath\LARGE\underline{円に内接$\cdot$外接する円}}\vspace{0.3zw}

\large
\bf\boldmath 問．中心が$\left(4,\;3\right)$で

\huge 
\hspace{0.1zw}
円$x^2+y^2=1$に\vspace{-0.1zw}\\
\hfill 接する円の方程式

\large
\hfill を求めよ．\\
\at(6.6cm,0.2cm){\small\color{bradorange}$\overset{\text{図形と方程式}}{\text{典型}}$}

\newpage

\at(0cm,0cm){\includegraphics[width=8cm,bb=0 0 1920 1080]{./youtube/thumbnails/templates/smart_background/図形と方程式.jpeg}}
{\color{orange}\bf\boldmath\LARGE\underline{円外から引いた接線}}\vspace{0.3zw}

\large
\bf\boldmath 問．点$\left(1,\;3\right)$から

\huge 
\hspace{0.1zw}
円$x^2+y^2=5$に\vspace{-0.1zw}\\
\hfill 引いた接線の方程式\vspace{0.2zw}

\large
\hfill を求めよ．\\

\at(6.6cm,0.2cm){\small\color{bradorange}$\overset{\text{図形と方程式}}{\text{典型}}$}

\newpage

\at(0cm,0cm){\includegraphics[width=8cm,bb=0 0 1920 1080]{./youtube/thumbnails/templates/smart_background/図形と方程式.jpeg}}
{\color{orange}\bf\boldmath\LARGE\underline{円と直線が接する条件}}\vspace{0.3zw}

\large
\bf\boldmath 問．
$\left\{\begin{array}{l} 円：x^2+y^2={10}\ \\直線：y=2x+m\ \end{array}\right.$が\\

\Huge
\vspace{-0.2zw}
\hspace{0.5zw}
接するとき，

\large
\vspace{0.4zw}
\hfill 定数$m$の値を求めよ．
\at(6.6cm,0.2cm){\small\color{bradorange}$\overset{\text{図形と方程式}}{\text{典型}}$}

\newpage

\at(0cm,0cm){\includegraphics[width=8cm,bb=0 0 1920 1080]{./youtube/thumbnails/templates/smart_background/図形と方程式.jpeg}}
{\color{orange}\bf\boldmath\Large\underline{円と直線が共有点をもつ条件}}\vspace{0.3zw}

\large
\bf\boldmath 問．
$\left\{\begin{array}{l} 円：x^2+y^2=8\ \\直線：y=x+m\ \end{array}\right.$が\\

\Huge
\vspace{-0.2zw}
\hspace{0.2zw}
共有点をもつとき，

\large
\vspace{0.4zw}
\hfill 定数$m$の値の範囲を求めよ．
\at(6.6cm,0.2cm){\small\color{bradorange}$\overset{\text{図形と方程式}}{\text{典型}}$}

\newpage

\at(0cm,0cm){\includegraphics[width=8cm,bb=0 0 1920 1080]{./youtube/thumbnails/templates/smart_background/図形と方程式.jpeg}}
{\color{orange}\bf\boldmath\LARGE\underline{円と直線の共有点}}\vspace{0.3zw}

\Large 
\bf\boldmath 問．円$x^2+y^2=5$と次の直線の共有点の座標を求めよ．\\
(1)  $y=x-1$\\
(2)  $y=2x+5$\\

\at(6.6cm,0.2cm){\small\color{bradorange}$\overset{\text{図形と方程式}}{\text{典型}}$}

\newpage

\at(0cm,0cm){\includegraphics[width=8cm,bb=0 0 1920 1080]{./youtube/thumbnails/templates/smart_background/図形と方程式.jpeg}}
{\color{orange}\bf\boldmath\LARGE\underline{$3$点を通る円の方程式}}\vspace{0.3zw}

\normalsize 
\bf\boldmath 問．$3$点$\text{A}\left(2,1\right),\text{B}\left(6,3\right),\text{C}\left(-1,2\right)$がある．

\normalsize
\vspace{0.3zw}
(1)  $3$点$\text{A},\;\text{B},\;\text{C}$を通る円の方程式を求めよ．

\LARGE
(2)  三角形$\text{ABC}$の外心の座標と，
外接円の半径を求めよ．\\

\at(6.6cm,0.2cm){\small\color{bradorange}$\overset{\text{図形と方程式}}{\text{典型}}$}

\newpage

\at(0cm,0cm){\includegraphics[width=8cm,bb=0 0 1920 1080]{./youtube/thumbnails/templates/smart_background/図形と方程式.jpeg}}
{\color{orange}\bf\boldmath\Large\underline{$2$直線の平行$\cdot$垂直条件}}\vspace{0.3zw}

\Large 
\bf\boldmath 問．$2$直線$\left\{\begin{array}{l}ax+2y=1,\\x+\left(a-1\right)y=3\end{array}\right.$が次の条件を満たすとき，定数$a$の値を求めよ．
(1)  平行 \hspace{0.2zw} (2)  垂直\\

\at(6.6cm,0.2cm){\small\color{bradorange}$\overset{\text{図形と方程式}}{\text{典型}}$}



\end{document}

